% =========================================================================
% CAPÍTULO 8
% =========================================================================
\chapter{Sistema de Alerta Temprana}

El sistema propuesto ingiere los boletines trimestrales de la CNBV, ejecuta el pipeline de limpieza e ingeniería de características, genera una probabilidad de \textit{Default} mediante el clasificador Stacking, y clasifica a cada institución en un nivel de riesgo:
\begin{itemize}
    \item \textbf{Riesgo bajo (verde):} probabilidad predicha $<0.30$.
    \item \textbf{Riesgo moderado (amarillo):} probabilidad entre 0.30 y 0.55.
    \item \textbf{Riesgo alto (rojo):} probabilidad $>0.55$.
\end{itemize}
Estos umbrales son una propuesta operativa basada en la distribución de probabilidades observada en la Figura \ref{fig:probdist}; su calibración final requeriría validación conjunta con el área de supervisión de la CNBV y no debe interpretarse como un estándar regulatorio ya adoptado.

En cuanto al flujo de decisión, ante una alerta roja, se recomienda: (i) revisión prioritaria por parte del supervisor de área; (ii) solicitud de plan de fortalecimiento de capital si el NICAP se aproxima a las categorías 2--3 definidas en el Capítulo 2; (iii) descomposición SHAP individual de la predicción para identificar los factores específicos que la originaron.

